\chapter*{The ideas behind the system}
\addcontentsline{toc}{chapter}{Introduction: The ideas behind the system}

Bijou's architecture starts with the conditions under which terminal software is used. A person may be navigating with a keyboard, inspecting a dense table, reading through assistive output or using a narrow viewport. An automated consumer may need the same information as structured facts. The system's design has to account for those encounters with the application.

The vision and doctrine chapters make the project's commitments explicit. Runtime truth gives priority to what the program actually does. Hexagonal boundaries separate application meaning from host services. The state-update-view model gives changes and effects identifiable owners. The layout and interaction principles connect rendered geometry to focus, selection and input routing.

\section*{How the principles fit together}

\subsection*{Meaning across modes}

Several themes recur. Graceful lowering asks how a rich surface preserves meaning in simpler output modes. Localization treats language and direction as inputs to layout and interaction. Inspectability asks what a person or an agent can learn about a surface without guessing from its appearance. The design language gives these concerns visible form through shared patterns, tokens and component families.

\subsection*{Principles under use}

DOGFOOD supplies a place to exercise those commitments in a real application. A principle becomes easier to assess when a reader can follow it into a shell, a component preview, a settings surface or a rendering test. This relationship between doctrine and a working surface explains the repeated emphasis on evidence throughout the later parts.

\section*{How to use these chapters}

Start with the vision, then the architecture and the invariants relevant to the work at hand. Treat an invariant as a constraint to investigate in the implementation. The prose states what the project is trying to make reliable; source code and focused checks establish how a particular revision meets that commitment.

\subsection*{Aspirations and transitions}

The strategy notes also preserve aspirations and transitions. Their wording belongs to their original context. When a later design changes a boundary or revises a policy, read the two records together and follow the linked integration evidence. This part provides the vocabulary for that reading; the design records show the vocabulary being put under pressure.
